%%%%%%%%%%%%%%%%%%%%%%%%%%%%%%%%%
% Abstract                      %
%%%%%%%%%%%%%%%%%%%%%%%%%%%%%%%%%

\section*{\hfill Abstract \hfill}

\addcontentsline{toc}{section}{Abstract}

\vspace{1.4cm}
\large
\noindent
The Tunisian banking sector is subject to growing regulatory pressure on data management, driven by the Banque Centrale de Tunisie Circulaire N\textdegree{}2025-08 and by the international standard BCBS 239, yet it lacks a standardized means of measuring and steering its data maturity. This project, conducted within the Data \& Analytics practice of EY Advisory Tunisia, addresses that gap by designing a data maturity assessment framework adapted to the banking sector and by developing an interactive tool that operationalizes it. The work follows a Design Science Research methodology for the design and validation of the framework, combined with Agile principles for the iterative development of the tool. A benchmark that screened fifteen international frameworks and retained five, namely DAMA-DMBOK, DCAM, CMMI-DMM, TDWI, and Gartner, together with a diagnostic survey of forty four banking professionals, informs the design. The resulting hybrid framework is structured around five dimensions, namely Governance, Data Quality, Architecture and Access, Analytics and Tools, and Skills and Culture, together with twelve sub-dimensions and forty seven evidence based indicators, of which thirteen are non-skippable regulatory indicators tied to the BCT Circulaire N\textdegree{}2025-08 and BCBS 239 requirements. A weighted and renormalized scoring model, protected by an evidence cap that prevents over-optimistic self-assessment, produces an auditable Global Maturity Index on a one to five scale. The framework is operationalized through DataPilot, an interactive steering tool that delivers automated scoring, visual dashboards, gap analysis, a regulatory compliance view, and a phased improvement roadmap, and that is delivered as a multi-tenant platform in which EY administers the EY master template, each bank receives its own copy, and four roles (EY Admin, Super Admin, Admin, Analyst) collaborate on solo and group assessments under database-enforced tenant isolation. Together, the framework and the tool provide Tunisian banks with a rigorous and actionable instrument to assess their data maturity and to support their digital transformation.

\vspace{0.6cm}
\noindent
\textbf{Keywords: Data maturity, Maturity assessment framework, Data governance, Banking sector, BCBS 239, Digital transformation, Multi-tenant platform, Web development, Software engineering, Design Science Research}
